\part*{ВВЕДЕНИЕ}
\addcontentsline{toc}{part}{\textbf{ВВЕДЕНИЕ}}

Когда компьютер работает в многозадачном режиме, на нем зачастую запускается сразу несколько процессов или потоков, претендующих на использование центрального процессора. Такая ситуация складывается в том случае, если в состоянии готовности одновременно находятся два или более процесса или потока. Если доступен только один центральный процессор, необходимо выбрать, какой из этих процессов будет выполняться следующим. Та часть операционной системы, на которую возложен этот выбор, называется \textbf{планировщиком}, а алгоритм, который он использует, называется \textbf{алгоритмом планирования}. В решении задачи планирования процессов ключевую роль играют приоритеты, время непрерывной работы, простоя и блокировки процессов. \cite{tannenbaum}

Linux является операционной системой разделения времени. Псевдопараллельное выполнение процессов обеспечивается за счёт квантования процессорного времени, позволяющего переключаться с одного процесса на другой за очень короткий промежуток времени.

%Та часть операционной системы, на которую возложен этот
%выбор, называется планировщиком, а алгоритм, который ею используется, называется
%алгоритмом планирования.
%
%Как и любая система с разделением времени, Linux достигает магического эффекта кажущегося
%одновременного выполнения нескольких процессов путем переключения с одного процесса на
%другой за очень короткий промежуток времени. Само переключение процессов обсуждалось в
%Глава 3; эта глава посвящена планированию, которое касается того, когда переключаться и какой процесс выбрать.
%когда переключаться и какой процесс выбрать.
%
%Для системы имеют большое значение приоритеты, время простоя и выполнения процессов. Эти значения используются для планирования. За планирование процессов в операционных системах отвечают планировщики задач, каждый из которых может использовать собственный уникальный алгоритм планирования процессов.

% Linux является операционной системой разделения времени. Классы планирования реального времени, блокировка памяти, разделяемая память и сигналы реального времени получили поддержку в Linux с самых первых дней. Очереди сообщений POSIX, часы и таймеры поддерживаются в ядре версии 2.6. Асинхронный ввод/вывод также поддерживается с самых первых дней, но эта реализация была полностью приведена в библиотеке языка Си пользовательского пространства. Linux версии 2.6 имеет поддержку AIO (Asynchronous I/O) в ядре. Библиотека языка C GNU и glibc также претерпели изменения для поддержки этих расширений реального времени. Для обеспечения лучшей поддержки в Linux POSIX.1b ядро и glibc работают вместе.

%Данная курсовая работа посвящена вопросам определения приоритетов, времени простоя выполнения процессов в операционной системе Linux, а также их мониторингу.

Целью данной работы является разработка загружаемого модуля ядра Linux для мониторинга изменения приоритетов, времени непрерывной работы, простоя и блокировки процессов.

